% Options for packages loaded elsewhere
\PassOptionsToPackage{unicode}{hyperref}
\PassOptionsToPackage{hyphens}{url}
\PassOptionsToPackage{dvipsnames,svgnames,x11names}{xcolor}
%
\documentclass[
  letterpaper,
  DIV=11,
  numbers=noendperiod]{scrartcl}

\usepackage{amsmath,amssymb}
\usepackage{iftex}
\ifPDFTeX
  \usepackage[T1]{fontenc}
  \usepackage[utf8]{inputenc}
  \usepackage{textcomp} % provide euro and other symbols
\else % if luatex or xetex
  \usepackage{unicode-math}
  \defaultfontfeatures{Scale=MatchLowercase}
  \defaultfontfeatures[\rmfamily]{Ligatures=TeX,Scale=1}
\fi
\usepackage{lmodern}
\ifPDFTeX\else  
    % xetex/luatex font selection
\fi
% Use upquote if available, for straight quotes in verbatim environments
\IfFileExists{upquote.sty}{\usepackage{upquote}}{}
\IfFileExists{microtype.sty}{% use microtype if available
  \usepackage[]{microtype}
  \UseMicrotypeSet[protrusion]{basicmath} % disable protrusion for tt fonts
}{}
\makeatletter
\@ifundefined{KOMAClassName}{% if non-KOMA class
  \IfFileExists{parskip.sty}{%
    \usepackage{parskip}
  }{% else
    \setlength{\parindent}{0pt}
    \setlength{\parskip}{6pt plus 2pt minus 1pt}}
}{% if KOMA class
  \KOMAoptions{parskip=half}}
\makeatother
\usepackage{xcolor}
\setlength{\emergencystretch}{3em} % prevent overfull lines
\setcounter{secnumdepth}{-\maxdimen} % remove section numbering
% Make \paragraph and \subparagraph free-standing
\makeatletter
\ifx\paragraph\undefined\else
  \let\oldparagraph\paragraph
  \renewcommand{\paragraph}{
    \@ifstar
      \xxxParagraphStar
      \xxxParagraphNoStar
  }
  \newcommand{\xxxParagraphStar}[1]{\oldparagraph*{#1}\mbox{}}
  \newcommand{\xxxParagraphNoStar}[1]{\oldparagraph{#1}\mbox{}}
\fi
\ifx\subparagraph\undefined\else
  \let\oldsubparagraph\subparagraph
  \renewcommand{\subparagraph}{
    \@ifstar
      \xxxSubParagraphStar
      \xxxSubParagraphNoStar
  }
  \newcommand{\xxxSubParagraphStar}[1]{\oldsubparagraph*{#1}\mbox{}}
  \newcommand{\xxxSubParagraphNoStar}[1]{\oldsubparagraph{#1}\mbox{}}
\fi
\makeatother


\providecommand{\tightlist}{%
  \setlength{\itemsep}{0pt}\setlength{\parskip}{0pt}}\usepackage{longtable,booktabs,array}
\usepackage{calc} % for calculating minipage widths
% Correct order of tables after \paragraph or \subparagraph
\usepackage{etoolbox}
\makeatletter
\patchcmd\longtable{\par}{\if@noskipsec\mbox{}\fi\par}{}{}
\makeatother
% Allow footnotes in longtable head/foot
\IfFileExists{footnotehyper.sty}{\usepackage{footnotehyper}}{\usepackage{footnote}}
\makesavenoteenv{longtable}
\usepackage{graphicx}
\makeatletter
\def\maxwidth{\ifdim\Gin@nat@width>\linewidth\linewidth\else\Gin@nat@width\fi}
\def\maxheight{\ifdim\Gin@nat@height>\textheight\textheight\else\Gin@nat@height\fi}
\makeatother
% Scale images if necessary, so that they will not overflow the page
% margins by default, and it is still possible to overwrite the defaults
% using explicit options in \includegraphics[width, height, ...]{}
\setkeys{Gin}{width=\maxwidth,height=\maxheight,keepaspectratio}
% Set default figure placement to htbp
\makeatletter
\def\fps@figure{htbp}
\makeatother

\KOMAoption{captions}{tableheading}
\makeatletter
\@ifpackageloaded{caption}{}{\usepackage{caption}}
\AtBeginDocument{%
\ifdefined\contentsname
  \renewcommand*\contentsname{Table of contents}
\else
  \newcommand\contentsname{Table of contents}
\fi
\ifdefined\listfigurename
  \renewcommand*\listfigurename{List of Figures}
\else
  \newcommand\listfigurename{List of Figures}
\fi
\ifdefined\listtablename
  \renewcommand*\listtablename{List of Tables}
\else
  \newcommand\listtablename{List of Tables}
\fi
\ifdefined\figurename
  \renewcommand*\figurename{Figure}
\else
  \newcommand\figurename{Figure}
\fi
\ifdefined\tablename
  \renewcommand*\tablename{Table}
\else
  \newcommand\tablename{Table}
\fi
}
\@ifpackageloaded{float}{}{\usepackage{float}}
\floatstyle{ruled}
\@ifundefined{c@chapter}{\newfloat{codelisting}{h}{lop}}{\newfloat{codelisting}{h}{lop}[chapter]}
\floatname{codelisting}{Listing}
\newcommand*\listoflistings{\listof{codelisting}{List of Listings}}
\makeatother
\makeatletter
\makeatother
\makeatletter
\@ifpackageloaded{caption}{}{\usepackage{caption}}
\@ifpackageloaded{subcaption}{}{\usepackage{subcaption}}
\makeatother

\ifLuaTeX
  \usepackage{selnolig}  % disable illegal ligatures
\fi
\usepackage{bookmark}

\IfFileExists{xurl.sty}{\usepackage{xurl}}{} % add URL line breaks if available
\urlstyle{same} % disable monospaced font for URLs
\hypersetup{
  pdftitle={Probability},
  pdfauthor={Stat 801a},
  colorlinks=true,
  linkcolor={blue},
  filecolor={Maroon},
  citecolor={Blue},
  urlcolor={Blue},
  pdfcreator={LaTeX via pandoc}}


\title{Probability}
\usepackage{etoolbox}
\makeatletter
\providecommand{\subtitle}[1]{% add subtitle to \maketitle
  \apptocmd{\@title}{\par {\large #1 \par}}{}{}
}
\makeatother
\subtitle{Random Variables, Probability Distributions}
\author{Stat 801a}
\date{}

\begin{document}
\maketitle


\subsection{Population versus Sample}\label{population-versus-sample}

{Population} -- The set of all objects under consideration. Objects in a
population are called {\emph{elements}}.

{Examples}

\begin{itemize}
\item
  Boxes of a particular brand of corn flakes.
\item
  Newly made light bulbs in a warehouse at a given time.
\end{itemize}

\subsection{Another Definition of
Population}\label{another-definition-of-population}

\begin{itemize}
\item
  All possible outcomes of an {\emph{experiment}} can also be values in
  a population.
\item
  Flipping a coin twice. Possible outcomes are:\\
  HH, HT, TH, TT. Could make the outcomes numeric using 3, 2, 1, 0 to be
  the elements corresponding to the above {events}.
\item
  If the experiment is repeated many times, the resulting values
  constitute a population: \(2,\ 0,\ 3,\ 1,\ 0,\ 2,\ \dots\)
\end{itemize}

\subsection{Random Sampling}\label{random-sampling}

\begin{itemize}
\item
  {Sample} -- A subset of the elements from a population.
\item
  {Simple Random Sample} of size \(n\) from a {finite population} is a
  subset selected in such a way that every subset of \(n\) elements is
  equally likely to be the one selected. (Abbreviation: SRS)
\item
  A random sample from a small finite population can be obtained using
  the {Random Numbers} - using code or random number table.

  \begin{itemize}
  \tightlist
  \item
    sample(1:100,3) or sample(3)
  \end{itemize}
\item
  Computer programs are used to draw samples from large populations.
\end{itemize}

\subsection{Probability, Random Variables, and Probability
Distributions}\label{probability-random-variables-and-probability-distributions}

\begin{itemize}
\item
  Assume a numeric finite population of \(N\) elements, and suppose
  \(n_1\) of the elements are the number \(k\).
\item
  Suppose a SRS of size 1 is taken from this population.
\item
  The {probability} that the number \(k\) is selected is \(n_1/N\).
\item
  That is, we define this probability to be the relative frequency of
  occurrence of \(k\) in the population.
\end{itemize}

\subsection{Probability: Example}\label{probability-example}

\begin{itemize}
\item
  {Example:} ~~\(N = 10\)
\item
  {Population}: 1, 1, 1, 3, 7, 7, 9, 9, 9, 9
\item
  Denote by \(P(k)\), {the probabilty of the number k being selected}
\item
  \(P(1) = 3/10\),~~\(P(3) = 1/10\),~~\(P(7) = 2/10\),\\
  \(P(9) = 4/10\),~~\(P(2) = 0/10\),~~\(P(1~\mbox{or}~3) = 4/10\),\\
  \(P(\mbox{not}~7) = 1-2/10\).
\end{itemize}

\subsection{Random Variables}\label{random-variables}

A {random variable} is a function defined on a numeric population. We
will often say that this function value is {the value assumed by the
random variable}.

\begin{itemize}
\item
  {Example}: \(N = 10\)
\item
  {Population}: 1, 1, 1, 3, 7, 7, 9, 9, 9, 9
\item
  Let \(X\) be a random variable defined on this population.
\item
  \(X\) can assume any of the values 1, 3, 7, or 9.
\end{itemize}

\subsection{Random Variables
(continued)}\label{random-variables-continued}

\begin{itemize}
\item
  A {Random variables} are different from ordinary variables because
  they assume values from the population with associated probabilities.
\item
  For e.g., we will say\\
  \strut ~~~\(P(X=1)=3/10\),~~~\(P(X=7)=2/10\),\\
  \strut ~~~\(P(X=2)=0/10\),~~~\(P(X=-6)=0/10\),\\
  \strut ~~~\(P(1 \leq X \leq 3)=4/10\),~~~\(P(1 \leq X \leq 9)=1\).
  etc.
\item
  A {probability distribution} specifies the probabilities associated
  with all possible values the random variable \(X\) can assume.\\
\end{itemize}

\subsection{The M \& M Example}\label{the-m-m-example}

{Example}:\\
If X represents the number of RED \(M\&Ms\) in a vending machine size
package, then perhaps {the probability distribution of \(X\)} is
something like the following:\\

\begin{longtable}[]{@{}cccccccc@{}}
\toprule\noalign{}
\(x\) & 5 & 6 & 7 & 8 & 9 & 10 & 11 \\
\midrule\noalign{}
\endhead
\bottomrule\noalign{}
\endlastfoot
\(P(X=x)\) & \(\frac{2}{28}\) & \(\frac{2}{28}\) & \(\frac{5}{28}\) &
\(\frac{7}{28}\) & \(\frac{5}{28}\) & \(\frac{4}{28}\) &
\(\frac{3}{28}\) \\
\end{longtable}

\subsection{\texorpdfstring{The M \(\&\) M Example
(continued)}{The M \textbackslash\& M Example (continued)}}\label{the-m-m-example-continued}

We read this table by finding a value of X, say 8, then reading the
probability below it:

\begin{itemize}
\tightlist
\item
  {The probability of observing 8 M\(\&\)Ms in a package is
  \(\frac{7}{28}\), meaning, if we open many many packages of M\(\&\)Ms,
  about 7 out of 28 will have exactly 8 M\&Ms within.}
\end{itemize}

This is an example of a {discrete random variable}.

\subsection{Discrete Random Variables}\label{discrete-random-variables}

\begin{itemize}
\item
  When the elements of a population consists of a countable discrete set
  of values the random variable defined on that population is called
  {discrete random variable}.
\item
  For discrete random variables the probability distribution can be
  displayed using a histogram or a table of values and associated
  probabilities as shown above.
\end{itemize}

\subsection{Properties of Probability Distributions of Discrete Random
Variable
X}\label{properties-of-probability-distributions-of-discrete-random-variable-x}

\begin{itemize}
\item
  \(0\leq P(X=x)\leq 1\) for all x
\item
  \(\sum_{x}P(X=x)=1\)
\end{itemize}

\subsection{The M \& M Probability
Distribution}\label{the-m-m-probability-distribution}

\includegraphics{part2_lectures_files/figure-pdf/unnamed-chunk-1-1.pdf}

\subsection{Continuous Random
Variables}\label{continuous-random-variables}

\begin{itemize}
\item
  When the population consists of an uncountable number of values in a
  line interval use integration to calculate proportions of the
  population that exist in any given {subinterval} of a continuum.
\item
  For a random variable defined on such a population we will deal with
  the probability that the random variable assumes a value in a given
  interval, not that it assumes a specific numeric value.
\end{itemize}

\subsection{Continuous Random Variables
(contd.)}\label{continuous-random-variables-contd.}

\begin{itemize}
\item
  When the population consists of a set of values in a continuous range
  the random variable defined on that population is called {continuous
  random variable}.
\item
  The probability distribution of a continuous random variable is a
  smooth curve.
\item
  The {area under the curve} between specified values of the random
  variable correspond to a probability.
\item
  The probability that a continuous random variable assumes a specific
  numeric value is zero.
\end{itemize}

\subsection{Continuous Random Variables
(contd.)}\label{continuous-random-variables-contd.-1}

\begin{itemize}
\item
  {Example:} A Continuous Random Variable \(X\) defined on
\item
  a {Population}: all numbers in the interval (0,2) occurring with
  various frequencies (that we won't specify here).
\item
  By Definition\\
  \(P(1/2 \leq X \leq 1/2) \equiv P(X = 1/2) = 0\)\\
  \(P(0 \leq X \leq 2) = 1,\ \ \ P(X < 0) = P(X \leq 0) = 0\)
\item
  We could possibly calculate the following probability\\
  \(P(1 \leq X \leq 2) = ?\)
\end{itemize}

\subsection{Binomial Random Variable}\label{binomial-random-variable}

\begin{itemize}
\item
  For some interesting populations that occur naturally, theoretical
  probability distributions and associated random variables, have been
  defined.
\item
  One such population features elements that are only one of two kinds
  (Yes or No; ~Good or Defective; ~Below or Above;, etc.).
\item
  Let us use {Success(S)} and {Failure(F)} to designate population
  elements, in general
\item
  Consider a typical population of this kind consisting of \(m\) {S}'s
  and \(N-m\) {F}'s, where \(N\) is the total number of elements in the
  population.
\end{itemize}

\subsection{Binomial Experiment}\label{binomial-experiment}

\begin{itemize}
\item
  A {binomial experiment} consists of taking \(n\) successive SRS, each
  of size 1, replacing elements between draws, and counting the number
  of {S}'s obtained.
\item
  The binomial experiment is conducted by sampling from the population
  of {S}'s and {F}'s.
\item
  The proportion of {S}'s in the population is \(\pi = m/N\) so the
  proportion of {F}'s is \(1 - \pi \equiv (N-m)/N\).
\item
  Let \(X_i\) be the random variable that assumes the value of the
  outcome of the \(i\)-th draw for each \(i=1,2,\ldots,n\).
\end{itemize}

\subsection{Binomial Distribution}\label{binomial-distribution}

\begin{itemize}
\item
  Let the random variable \(X_i\) assume the value 1 if the outcome of
  the \(i\)-th draw is {S} or the value 0 if the outcome is {F}.\\
\item
  Thus we imagine \(n\) random variables, \(X_1, X_2, X_3, \ldots, X_n\)
  each assuming a value {S} with probability \(\pi\), or {F} with
  probability \(1 - \pi\).\\
\item
  That is, the probability distribution of \emph{each} of the \(n\)
  random variables \(X_1, X_2, X_3, \ldots, X_n\) is given by:\\

  \begin{longtable}[]{@{}ccc@{}}
  \toprule\noalign{}
  \(x\) & 0 & 1 \\
  \midrule\noalign{}
  \endhead
  \bottomrule\noalign{}
  \endlastfoot
  \(P(X=x)\) & \(1-\pi\) & \(\pi\) \\
  \end{longtable}
\end{itemize}

\subsection{Binomial Distribution
(contd.)}\label{binomial-distribution-contd.}

\begin{itemize}
\item
  Now define a new random variable \(Y\) so that
  \[Y = X_1 + X_2 + X_3 + \cdots + X_n \,\]
\item
  \(Y\) represents the number of {Success(S)}'s in the set of values
  assumed by \(X_1,\ X_2, \ldots,X_n\).\\
\item
  This is because each \(X_i=1\) when an {S} is drawn, and \(X_i=0\)
  otherwise, we can see that \(Y\) represents the total number of {S}'s
  in the experiment.\\
\item
  \(Y\) is said to be a {binomial random variable}.
\end{itemize}

\subsection{Binomial Distribution
(contd.)}\label{binomial-distribution-contd.-1}

\begin{itemize}
\item
  That is, a {binomial random variable} assumes the value of the total
  number of successful outcomes in a binomial experiment.\\
\item
  The random variable \(Y\) can assume any one of the values 0, 1, 2,
  \(\ldots, n\) with a probability that will depend on \(n\), the number
  of samples taken and \(\pi\).\\
\item
  \(Y\) is said to be the Binomial (\(n,\pi\)) random variable.\\
\item
  A population on which \(Y\) is defined consists of the unique elements
  \(0,1,2,\ldots,n\).\\
\item
  What is the proportion of each element in the population? i.e., What
  is the probability that \(Y\) assumes each value?
\end{itemize}

\subsection{Binomial Distribution
(contd.)}\label{binomial-distribution-contd.-2}

Theoretically, for any finite \(n\), the probability that a\\
{binomial random variable} \(Y\) takes on values 0, 1, 2, \(\dots, n\)\\
is given by\\
\[\begin{eqnarray*}
P(Y=y) & = & \frac{n!}{y!(n-y)!}\pi^y\,(1 - \pi)^{n-y}, \\
\mbox{where} & & \\
        y & = & 0, 1, 2, \ldots, n \mbox {(number of successes in n trials)}\\
\pi & = &  \mbox{ probability of success in a single trial}\\
1-\pi & = &  \mbox{ probability of failure in a single trial}\\
n! & = & n(n-1)(n-2) \ldots (3)(2)(1)
\end{eqnarray*}\]

\subsection{Binomial Distribution
(contd.)}\label{binomial-distribution-contd.-3}

\begin{itemize}
\item
  In applications, we won't know \(\pi\) and our objective will usually
  be to sample the population in order to {estimate} \(\pi\), and get an
  idea about the {error} in that estimate.
\item
  What is the {mean} \(\mu\) and {variance} \(\sigma^2\) (or the
  standard deviation) \(\sigma\) of a binomial random variable?
\item
  In general, for a {discrete random variable}, these are defined is as
  follows.
\end{itemize}

\subsection{Computing the Mean and Variance of a Discrete Random
Variable}\label{computing-the-mean-and-variance-of-a-discrete-random-variable}

\begin{itemize}
\item
  The {mean} of a discrete random variable \(X\) is:
  \[E(X) = \mu = \sum_x\,x \cdot P(X=x)\]
\item
  The {variance} of a discrete random variable \(X\) is:
  \[V(X) = \sigma^2= \sum_x\,(x - E(X))^2\cdot \ P(X=x)\]
\item
  The {standard deviation} of a random variable \(X\) is
  \[\sqrt{V(X)}=\sigma\]
\end{itemize}

\subsection{Computing the Mean and Variance of a Binomial Random
Variable}\label{computing-the-mean-and-variance-of-a-binomial-random-variable}

\begin{itemize}
\item
  Using the above formula, the {mean} of a Binomial random variable
  ~\(Y\), is shown to be \(E(Y)=\mu = n \, \pi\)
\item
  The {variance} is shown to be \(V(Y) =\sigma^2= n\,\pi(1 - \pi)\).
\item
  Thus the {standard deviation} of a Binomial random variable is
  \(\sigma= \sqrt{n\,\pi(1 - \pi)}\)
\end{itemize}

\subsection{Binomial Distribution:
Example}\label{binomial-distribution-example}

Consider a Binomial population with \(n=3,\ \pi = \frac13\),\\
i.e., \(Y\ \sim\ \mbox{Bin}(3,\,\frac13)\), The corresponding Binomial
distribution is\\
\[\begin{eqnarray*}
  P(0) &=& \frac{3!}{0!3!} \left(\frac13 \right)^0 \left( \frac23
           \right)^3 = 8/27\\[.1in]
  P(1) &=&  \frac{3!}{1!2!} \left(\frac13 \right)^1 \left( \frac23
           \right)^2 = 3 \cdot \left( \frac13
           \right)\, \left(\frac49 \right) = 12/27 \\[.1in]    
  P(2) &=&  \frac{3!}{2!1!} \left(\frac13 \right)^2 \left( \frac23
           \right) = 3 \cdot \left( \frac19
           \right)\, \left(\frac23 \right) = 6/27 \\[.2in]
  P(3) &=& 1 - P(0) - P(1) - P(2) = 1/27
\end{eqnarray*}\]

\subsection{Binomial Distribution: Example
(contd.)}\label{binomial-distribution-example-contd.}

The mean and variance of a Bin(\(3,\frac{1}{3}\)) population is:

\begin{itemize}
\item
  Mean \(= n \, \pi =  3 \times\frac{1}{3} = 1\)
\item
  Variance
  \(=  n\,\pi(1 - \pi) =  3 \times\frac{1}{3}\times\frac{2}{3} =
  \frac{2}{3}\)
\item
  Standard Deviation \(= \sqrt{.66667} = .8165\)
\end{itemize}

\subsection{Using the Binomial
Distribution}\label{using-the-binomial-distribution}

\begin{itemize}
\item
  So how can this Binomial probability distribution benefit us?
\item
  When sampling situation fits the requirements of a Binomial
  experiment, (at least approximately) we can view the results we obtain
  as the values assumed by a binomial random variable.
\item
  This permits us to find probabilities of various possible outcomes,
  and use these and other properties of a binomial random variable to
  make inferences.
\item
  To assume the {Binomial Model}, we need to determine whether the
  values of the population elements are obtained as a result of a
  {Binomial Experiment}.
\end{itemize}

\subsection{The Binomial Model}\label{the-binomial-model}

\begin{enumerate}
\def\labelenumi{\arabic{enumi}.}
\item
  \(n\) identical trials are made.
\item
  Each trial results in either {Success(\(S\))} or {Failure(\(F\))}.
  (only two possible outcomes).
\item
  The probability of success (\(\pi\)) remains the same from trial to
  trial.
\item
  The trials are independent.
\item
  Interest is in the number of successes obtained in the \(n\) trials.
  The number is viewed as being the value assumed by a Binomial
  (\(n,\ \pi\)) random variable \(Y\).
\end{enumerate}

\subsection{The Binomial Model}\label{the-binomial-model-1}

\begin{itemize}
\item
  Requirement Number 3 in the definition of a binomial experiment is
  usually the most difficult to satisfy in practical applications.
\item
  As trials are made, if selected elements are not replaced in the
  sampled population, \(\pi\) changes in value.
\item
  If the population is very large, the change is very small and can be
  ignored.
\end{itemize}

\subsection{Binomial Model: Are requirements satisfied? - Example
1}\label{binomial-model-are-requirements-satisfied---example-1}

An article in the March 5, 1998, issue of The New England Journal of
Medicine discussed a large outbreak of tuberculosis. One person, called
the index patient, was diagnosed with tuberculosis in 1995. The 232
co-workers of the index patient were given a tuberculin screening test.
The number of co-workers recording a positive reading on the test was
the random variable of interest. Did this study satisfy the properties
of a binomial experiment?

\subsection{Binomial Model: Are requirements satisfied? - Example
1}\label{binomial-model-are-requirements-satisfied---example-1-1}

\begin{enumerate}
\def\labelenumi{\arabic{enumi}.}
\item
  \(n\) identical trials? Yes, \(n=232\) workers
\item
  Population of only \(S\)'s and \(F\)'s? Yes.
\item
  Is probability of success same in each trial? Yes if we assume the
  co-workers had the same risk factors as the index patient.
\item
  Independent trials? Yes, the outcome of one screening was unaffected
  by the outcome of the other tests.
\item
  Was the random variable of interest to the experimenter the number of
  successes \(y\) in the 232 screening tests? Yes, the number of
  co-workers who received a positive result.
\end{enumerate}

\subsection{Binomial Model: Are requirements satisfied? - Example
2}\label{binomial-model-are-requirements-satisfied---example-2}

An economist interviews 75 students in a class of 100 to estimate the
proportion of student who expect to obtain a C or better in the course.
Is this a binomial experiment?

\subsection{Binomial Model: Are requirements satisfied? - Example
2}\label{binomial-model-are-requirements-satisfied---example-2-1}

\begin{itemize}
\item
  Population of only \(S\)'s and \(F\)'s? Yes.
\item
  Is probability of success same in each trial? This may change from
  trial to trial as we may not assume sampling with replacement.
\end{itemize}

\subsection{Binomial Example}\label{binomial-example}

A labor union's examining board for the selection of apprentices has a
record of admitting 70\% of all applicants who satisfy a basic set of
requirements. Five members of a minority group, all satisfying the
requirements, recently came before the board, and four of the five were
rejected. Find the probability that 1 or fewer would be accepted if the
admission rate is really 0.70.

\subsection{Binomial Example - cont.}\label{binomial-example---cont.}

\begin{itemize}
\item
  Board decisions are viewed as random draws from a population
  consisting of \(70\%\) Accepts (S) and \(30\%\) Rejects (F).
\item
  We can assume a large population so that the probability of being
  accepted remains constant across draws.
\item
  Then the Binomial (5,0.7) distribution is applicable to the number of
  Accepts in \(n=5\) trials with \(\pi=0.7\).
\end{itemize}

\subsection{Binomial Example - cont.}\label{binomial-example---cont.-1}

\begin{itemize}
\item
  Let \(X\) be the random variable representing the number of Accepts,
  and assume that \(X\) has a Binomial (5, 0.7) distribution. Then
  P(X=0) + P(X=1) =\\
  \({5}\choose{0}\) \((0.7)^0(0.3)^5+\) \({5}\choose{1}\)
  \((0.7)(0.3)^4=0.03078\)
\item
  The probability that the board will accept fewer than 2 applicants
  from a group of 5 is 0.03078 (which is a small probability) {if indeed
  the true acceptance rate is 0.7.}
\end{itemize}

\subsection{Binomial Example - cont.}\label{binomial-example---cont.-2}

\begin{itemize}
\item
  But in fact the board DID accept fewer than 2 of the group of 5 so
  .......\ldots.
\item
  {Conclusion:} Either the board was operating to admit at its
  historical 70\% rate and we have simply witnessed an extremely rare
  event, {or}
\item
  the board has changed policy and is currently admitting at a rate less
  than 70\%.
\item
  One would tend to favor the latter conclusion.
\end{itemize}

\subsection{Normal Random Variable}\label{normal-random-variable}

\begin{itemize}
\item
  The {Normal probability distribution} is an example of a continuous
  distribution.
\item
  The associated random variable is defined in the interval
  (-\(\infty,\infty)\).
\item
  To model the probability distribution of many populations we often use
  the Normal distribution.
\item
  Most population {distributions} are bell-shaped and can be modeled by
  the Normal distribution.
\item
  We say that the associated {random variable} \(X\), say,\\
\item
  Here \(\mu\) and \(\sigma\), are {parameters} of the population.
\end{itemize}

\subsection{Properties of the Normal
Distribution}\label{properties-of-the-normal-distribution}

\begin{itemize}
\item
  Probability function (or curve) of a Normal random variable \(X\) is
  \[\begin{equation*}
    f(x) = \frac{1}{\sigma \sqrt{2\pi}} \, e^{- \frac{1}{2} \left(\frac{x-\mu}{\sigma}\right)^2}
  \end{equation*}\] where
\item
  \(E(X)= \mu\) = mean : measures the {center of the distribution}.
\item
  \(V(X)= \sigma^2\) = variance : {measures the spread}.
\item
  \(\sqrt{V(X)}= \sigma\) = standard deviation : {measures the spread in
  same units as the mean}.
\item
  \(\pi\) = 3.1415\ldots{} \(e\) = 2.71828\ldots{}
\end{itemize}

\subsection{Properties of the Normal Distribution
(contd.)}\label{properties-of-the-normal-distribution-contd.}

\begin{itemize}
\item
  The Normal distribution is symmetric about the mean \(\mu\).
\item
  There are infinitely many Normal curves, one for each pair of values
  of \(\mu\) and \(\sigma\).
\item
  We need to compute the areas under the curve for many of these for
  statistical inference.
\item
  Mathematically speaking, proportions of a Normal\((\mu, \sigma^2)\)
  population which fall in various intervals (say (\(a,b\)) is one such
  interval) are obtained by evaluating the following integral (for
  specified values of \(\mu, \sigma^2\), \(a\) and \(b\)).
  \[\begin{equation*}
    P(a \leq X \leq b) = \frac{1}{\sqrt{2\pi}\sigma}
    \int_a^b \, e^{-\frac{1}{2}\left(\frac{x-\mu}{\sigma}\right)^2} dx
  \end{equation*}\]
\end{itemize}

\subsection{The Standard Normal
Distribution}\label{the-standard-normal-distribution}

\begin{itemize}
\item
  Table 1 lists the computed areas for the Normal distribution with
  \(\mu = 0\), \(\sigma = 1\).
\item
  We call this the {standard normal distribution} and denote it by
  \(N(0,1)\).~
\item
  The relationship between a standard normal random variable \(Z\) and a
  normal random variable \(X \sim N(\mu, \sigma^2)\) is
  \[Z = \frac{X-\mu}{\sigma}.\]
\end{itemize}

\subsection{The Standard Normal Distribution
(contd.)}\label{the-standard-normal-distribution-contd.}

\begin{itemize}
\item
  Thus (since \(\sigma > 0\)) for any pair of numbers \(a \leq b\),
  \[\begin{eqnarray*}
       P(a \leq X \leq b) & = & P(\frac{a-\mu}{\sigma} \leq
  \frac{X-\mu}{\sigma} \leq \frac{b-\mu}{\sigma})\\
             & = & P(\frac{a-\mu}{\sigma} \leq
  Z \leq \frac{b-\mu}{\sigma})
  \end{eqnarray*}\]
\item
  Thus we can compute probabilities for \(X\) for any member of the
  normal family using the probabilities associated with \(Z\) which are
  available from the Standard Normal Probabilities table.
\end{itemize}

\subsection{The Standard Normal Distribution
(contd.)}\label{the-standard-normal-distribution-contd.-1}

\begin{itemize}
\item
  The \(p\) quantiles of a Normal (\(\mu,\sigma^2\)) population are real
  numbers (use \(x_p\) to denote the \(p\) quantile) such that
  \[p = \int_{-\infty}^{x_p} \;
  \frac{1}{\sqrt{2\pi}\sigma} \, \,e^{-\frac12(\frac{x-\mu}{\sigma})^2} \; dx\]
\item
  If you can find a value \(x\) that satisfies the above equation then
  that value will be the \(p^{th}\) quantile of the \(N(\mu,\sigma^2)\)
  population i.e., \(x = Q(p)\).
\item
  The {standard normal quantiles} are the \(z\)-values found from the
  Standard Normal Probabilities table for specified \(p\) values that
  are in the body of the table.
\end{itemize}

\subsection{Examples: Using the Standard Normal Probabilities
table}\label{examples-using-the-standard-normal-probabilities-table}

\begin{enumerate}
\def\labelenumi{\arabic{enumi}.}
\item
  \(P(Z \geq 0) = P(Z \leq 0) = 0.50\)
\item
  \(P(0 \leq Z \leq 0.12) \equiv P(0 < Z < 0.12) \hspace*{1.5in}=P(Z<0.12)-P(Z<0)=0.5478-0.5 = 0.0478\)
\item
  \(P(-0.12 \leq Z \leq 0) =P(Z<0)-P(Z<-0.12)
   = 0.5-.4522=0.0478\)
\item
  \(P(-0.12 \leq Z \leq 0.12) = 2 \cdot P(0 \leq Z \leq 0.12) = 0.0956\)
\item
  \(P(0 \leq Z \leq 3) = 0.9987-0.5=0.4987\)
\item
  \(P(1 \leq Z \leq 3 ) = P(Z \leq 3)- P(Z \leq 1) = 0.9987 - 0.8413 \hspace*{1in}=0.1574\)
\end{enumerate}

\subsection{Examples: Using the Standard Normal Probabilities
table}\label{examples-using-the-standard-normal-probabilities-table-1}

\begin{enumerate}
\def\labelenumi{\arabic{enumi}.}
\item
  Find \(z\) such that \(P(0 \leq Z \leq z) = 0.4846\) Equivalent to
  finding \(z\) such that \(P(Z \leq z) = 0.9846\) which gives
  \(z = 2.16\)
\item
  Find \(z\) such that \(P(z \leq Z \leq 0) = 0.2257\)\\
  Equivalent to finding \(z\) such that \(P(Z \leq z) = 0.2743\) which
  gives \(z = -0.60\)
\item
  Find \(z\) such that \(P(Z \geq z) = 0.011\)\\
  Equivalent to finding \(z\) such that \(P(Z \leq z) = 0.9890\) which
  gives \(z = 2.29\)
\end{enumerate}

\subsection{\texorpdfstring{Examples for
\(X \sim N(\mu, \sigma^2)\)}{Examples for X \textbackslash sim N(\textbackslash mu, \textbackslash sigma\^{}2)}}\label{examples-for-x-sim-nmu-sigma2}

\begin{itemize}
\item
  When \(\mu = 1\), \(\sigma = 2\), calculate \(P(0 \leq X \leq 1)\) =
  \(P(\frac{0-1}{2} \leq Z \leq \frac{1-1}{2})= P(-0.5 \leq Z \leq 0) = 0.5-0.3085=0.1915\)
\item
  When \(\mu = 3\), \(\sigma=4\),calculate \(P(2 \leq X \leq 4)\) =
  \(P(\frac{2-3}{4} \leq Z \leq \frac{4-3}{4})\) =
  \(P(- 0.25 \leq Z \leq 0.25)\) \(= 0.5987-0.4013=0.1974\)
\item
  Find the 0.25 quantile of the normal population having \(\mu = 3\),
  \(\sigma = 4\). Find x so that \(P(X \leq x) = 0.25,\) or
  \(P(Z \leq \frac{x-3}{4}) = 0.25\)
  \(\frac{x-3}{4} \doteq -0.67 \Rightarrow x \doteq 3+ 4 \times (-0.67)= 0.32\)
\end{itemize}

\subsection{\texorpdfstring{Sampling Distribution of Sample Mean
\(\bar{Y}\).}{Sampling Distribution of Sample Mean \textbackslash bar\{Y\}.}}\label{sampling-distribution-of-sample-mean-bary.}

\begin{itemize}
\item
  The outcome of drawing a random sample of size \(n\) can be
  characterized as the value realized by each of \(n\) {random
  variables} \(Y_1, Y_2, \ldots, Y_n\), each defined independently on
  the population.
\item
  Thus the {Sample Mean} \[\overline{Y} = \frac1n \sum_{i=1}^n Y_i\] (or
  any other function of \(Y_1, Y_2, \ldots, Y_n\)) is a {random
  variable}.
\end{itemize}

\subsection{\texorpdfstring{Sampling Distribution of Sample Mean
\(\bar{Y}\).}{Sampling Distribution of Sample Mean \textbackslash bar\{Y\}.}}\label{sampling-distribution-of-sample-mean-bary.-1}

\begin{itemize}
\item
  The probability distribution of \(\overline{Y}\) is said to be the
  {Sampling Distribution} of \(\overline{Y}\).
\item
  It is the {derived population} of means of all possible samples of
  size \(n\) taken from the sampled population.
\item
  Properties of \(\overline{Y}\) include \(E(\overline{Y}) = \mu\), and
  \(V(\overline{Y}) = \sigma^2/n,\) where \(\mu\) and \(\sigma^2\) are
  the population mean and variance of the original sampled population.
\item
  The standard deviation of \(\overline{Y}\) is called {Standard Error
  of the Mean} and is written \(\sigma_{\overline{Y}}\).
\end{itemize}

\subsection{The Central Limit Theorem}\label{the-central-limit-theorem}

\begin{itemize}
\item
  The following theorem shows why the normal family of distributions is
  so important.
\item
  Assume sampling is from some finite or infinite population with
  parameters \((\mu,\,\sigma^2)\). Let \(Y_1, Y_2, \ldots, Y_n\) denote
  the random variables corresponding to a sample of size \(n\)
\item
  Let \(\overline{Y} = \frac1n \displaystyle\sum_{i=1}^n Y_i\).
\item
  {Central Limit Theorem (CLT)}\\
  When sample size \(n\) is large, the random variable \(\bar{Y}\) is
  approximately normally distributed with mean \(\mu\) and variance
  \(\sigma^2/n\). (i.e., \(\overline{Y} \approx N (\mu, \sigma^2/n\)))
\end{itemize}

\subsection{Using the Central Limit
Theorem}\label{using-the-central-limit-theorem}

\begin{itemize}
\item
  The approximation can be very close for \(n\) as small as 10 if the
  sampled population elements are symmetrically distributed around
  \(\mu\).
\item
  For mildly skewed sampled populations, \(n=30\) can give a close
  approximation. For higher skewed populations larger \(n\) will be
  needed for a better approximation.
\item
  So, the CLT says that regardless of the distribution of the sampled
  population, the derived population of the sample means is
  approximately a normal population and the approximation gets better as
  \(n\) gets larger.
\end{itemize}

\subsection{Central Limit Theorem:
Example}\label{central-limit-theorem-example}

Estimating average weight of sheep in a large herd. Random sample of
size \(n=50\). Assume \(\sigma^2=4\). but \(\mu\), population mean
animal weight in the herd is unknown.

\begin{itemize}
\item
  What is the approximate probability distribution of the sample mean
  random variable?
\item
  If \(\mu=40\) what is the probability that the sample mean will be
  within 0.1 pounds of \(\mu\)?
\end{itemize}

Answer:

\begin{itemize}
\tightlist
\item
  By the CLT, \(\overline{Y} \approx N (\mu, 4/50)\)
\end{itemize}

\subsection{Example continued}\label{example-continued}

\[\begin{eqnarray*}
  {P(39.9 \leq \bar{Y} \leq 40.1)  \doteq }\\[.1in]
  & & P\left(\frac{39.9-40}{\sqrt{4/50}}\right. 
   \left.\leq z \leq \frac{40.1-40}{\sqrt{4/50}}\right) \\[.2in]
     & = & P\left(\frac{-\sqrt{50}}{20} \leq z \leq 
         \frac{\sqrt{50}}{20}\right) \\[.3in]
     & = & P\left(-0.35 \leq z \leq 0.35\right) \\[.3in]
        & = & 0.6368-0.3632=0.2736
  \end{eqnarray*}\]

\subsection{Exercise}\label{exercise}

A random sample of 16 measurements is drawn from a population with a
mean of 60 and a standard deviation of 5. Describe the sampling
distribution of \(\bar{y}\), the sample mean. Within what
interval,symmetric about 60, would you expect \(\bar{y}\) to lie
approximately 95\% of the time?\\
Solution to Exercise:

\begin{itemize}
\item
  \(\bar{Y} \approx N(60, \frac{25}{16})\)
\item
  \(Z = \frac{\bar{Y}-60}{5/4} \approx N(0,1)\)
\item
  We know \(P(-z_{\frac{0.05}{2}} \leq Z \leq z_{\frac{0.05}{2}}) =
           0.95\) and \(z_{0.025} = 1.96\) from the Standard Normal
  Probabilities Table
\end{itemize}

\subsection{Solution to Exercise
(cont.)}\label{solution-to-exercise-cont.}

\begin{itemize}
\item
  Thus \(P(-1.96 \leq Z \leq 1.96) =.95\) which implies that a value of
  \(Z\) is {expected} to lie in the interval \([-1.96, \, 1.96]\)
  approximately 95\% of the time
\item
  \(-1.96 \leq \frac{\bar{y}-60}{5/4} \leq 1.96\)\\
  \(60 - 1.96 \times \frac54 \leq \bar{y} \leq 60 +
               1.96 \times \frac54\) giving
  \(57.55 \leq \bar{y} \leq 62.45\)\\
  Thus \([57.55, \, 62.45]\) is the interval in which \(\bar{y}\), the
  is {expected} to lie approximately 95\% of the time
\end{itemize}

\subsection{Exercise 2}\label{exercise-2}

Psychomotor retardation scores for a large group of manic-depressive
patients were found to be approximately normal with a mean of 930 and a
standard deviation of 130.

\begin{itemize}
\item
  What fraction of the patients scored between 800 and 1,100?
\item
  Less than 800?
\item
  Greater than 1200?
\item
  Find the 90th percentile for the distribution of manic-depressive
  scores.
\item
  Find the interquartile range.
\end{itemize}

\subsection{Solution to Exercise 2:}\label{solution-to-exercise-2}

d.) \(P(Z < a) = 0.90\) From the Standard Normal Probabilities Table, a
= 1.28\\
Set \(\frac{y - 930}{130} = 1.28\) and solving for \(y\) gives
\(y = 930 + 166.4\)\\
Thus \(= 1096.4\) is the 90th percentile; or the 0.9 quantile.\\
e.) \(\begin{array}{l|l}
            P(Z < a) = 0.75 & P(Z < a) = 0.25\\
            \mbox{gives } a = 0.675 & \mbox{gives } a = -0.675\\[.1in]
            \mbox{Set } \frac{y-930}{130} = 0.675 & \mbox{Set } \frac{y-930}{130} = -0.675\\[.1in]
            \mbox{Thus } Q(0.75) = 930. + 87.75 & \mbox{Thus } Q(0.25)= 930 -87.75  \\
            = 1017.25 & = 842.25
            \end{array}\)\\
Thus~ IQR \(= 1017.25- 842.25 = 175.50\)

\subsection{Another Exercise}\label{another-exercise}

The breaking strengths for 1-foot-square samples of a particular
synthetic fabric are approximately normally distributed with a mean of
2,250 pounds per square inch (psi) and a standard deviation of 10.2
psi.\\

\begin{itemize}
\item
  Find the probability of selecting a 1-foot-square sample of material
  at random that on testing would have a breaking strength in excess of
  2,265 psi.
\item
  Describe the sampling distribution for \(\bar{y}\) based on random
  samples of 15 one-foot sections.
\end{itemize}

\subsection{Another Exercise -
Solution}\label{another-exercise---solution}

\begin{itemize}
\item
  \(Y \sim N(2250, (10.2)^2)\)\\
  \[\begin{eqnarray*}
              P(Y > 2265) & = & P\left(Z > \frac{2265-2250}{10.2}\right)\\[.2in]
              & = & P(Z > 1.47) =  0.07078
  \end{eqnarray*}\]
\item
  \(\overline{Y} \approx N\left(2250, \frac{(10.2)^2}{15}\right)\)\\
  i.e., \(\overline{Y}\) is approximately distributed as Normal with
  mean=\(\mu_{\bar{y}}=2250\) and standard
  deviation=\(\sigma_{\bar{y}} = 2.63\)
\end{itemize}

\subsection{Poisson Distribution}\label{poisson-distribution}

\begin{itemize}
\item
  Poisson distribution is used to model the number of occurrences of
  ``rare'' events in certain intervals of time and space.
\item
  Certain conditions must be met regarding the occurrence of these
  events and the definition of the period of time or region of space in
  which they occur for a Poisson random variable to represent the number
  of events.
\item
  A Poisson random variable is a {discrete} random variable as it
  represents a \emph{count}.
\end{itemize}

\subsection{Poisson Distribution -
Examples}\label{poisson-distribution---examples}

\begin{itemize}
\item
  \(X\) = \# of alpha particles emitted from a polonium bar in an 8
  minute period\\
\item
  \(Y\) = \# of flaws on a standard size piece of manufactured product
  (e.g., 100m coaxial cable, 100 sq.meter plastic sheeting)\\
\item
  \(Z\) = \# of hits on a web page in a 24h period
\end{itemize}

{Definition:} The Poisson probability mass function is defined as:\\
\(p(x) = \frac{e^{-\lambda}\lambda^{x}}{x!} \hspace{1cm} \rm{ for }\ x = 0,1,2,3,\ldots\)\\
\(\lambda\) is called the {\emph{mean parameter}}, and represents the
\emph{expected number} of events occuring in the interval of interest.\\

\subsection{Computing Mean and Variance of a Poisson Random
Variable}\label{computing-mean-and-variance-of-a-poisson-random-variable}

\begin{itemize}
\item
  Expected Value of \(X \sim Poi({\lambda})\) is:

  \begin{itemize}
  \tightlist
  \item
    \(E[X]= \sum_{x=0}^{\infty} x\frac{e^{-\lambda}\lambda^{x}}{x!} = e^{-\lambda}\sum\limits_{x=0}^\infty\frac{\lambda^x}{(x-1)!}=e^{-\lambda}\lambda\sum\limits_{x=1}^\infty x\frac{\lambda^{x-1}}{(x-1)!}= e^{-\lambda}\lambda\sum\limits_{x=0}^\infty x\frac{\lambda^x}{x!}=\lambda\)
  \end{itemize}
\item
  Variance of \(X \sim Poi({\lambda})\) is:

  \begin{itemize}
  \tightlist
  \item
    \(\rm{Var}[X] =\sum_{x=0}^{\infty} (x-\lambda)^2 \{\frac{e^{-\lambda}\lambda^{x}}{x!}\} =\ldots = \lambda\)
  \end{itemize}
\end{itemize}

\subsection{Poisson distribution:
Example}\label{poisson-distribution-example}

A manufacturer of chips produces 1\% defectives. What is the probability
that in a box of 100 chips no defectives are found?\\
{Solution:} A defective chip can be considered to be a rare event, since
\(p\) is small (\(p = 0.01\)). So, model \(X\) as {Poisson} random
variable.\\
{We need to obtain a value for \(\lambda\)! Look at the expected
value!}\\
Note that we expect \(100 \cdot 0.01\) = 1 chip out of the box to be
defective.\\
{We know that the expected value of \(X\) is \(\lambda\). In this
example, therefore, we take \(\lambda = 1\).}\\
Then \(P(X=0) = \frac{e^{-1}1^{0}}{0!} = 0.3679.\)

\subsection{Poisson distribution: Example
(cont.)}\label{poisson-distribution-example-cont.}

{For the same product, what is the probability that, in a box of 200
chips, no defectives are found?}\\
{Solution:} The number of defectives in a box of 200 chips \(Y\) has a
Poisson distribution with parameter \(\lambda=2\) (since the expected
number of defectives in 200 chips is 2.)\\
Then, we have \(P(Y=0)= .1353\)\\
Note that \(X\) and \(Y\) are random variables with different Poisson
distributions because the events they represent occur in different
regions (two different boxes here).\\

\subsection{Poisson to approximate
Binomial}\label{poisson-to-approximate-binomial}

{Result (not a theorem):} For large \(n\), the Binomial distribution can
be approximated by the Poisson distribution, where \(\lambda\) is taken
as \(n\pi\):
\[{n \choose k} \pi^{k}(1-\pi)^{n-k} \approx e^{-n\pi}\frac{(n)^{k}}{k!}\]
{Rule of thumb:} use Poisson approximation if
\(n \ge 100,\ \pi \le 0.01\) and \(n\pi \leq 20\)\\




\end{document}
